% 2013 湖北卷理科第 8 题：四段叠放几何体的三视图（TikZ 矢量重绘）。
% 正视图自上而下为：圆台（上径 4、下径 2、高 1）、圆柱（径 2、高 2）、
% 正四棱柱（2×2、高 2）、正四棱台（上 2、下 4、高 1）。图中尺寸单位与题设一致。
\begin{tikzpicture}[scale=0.52, line width=0.5pt, font=\small, >=Stealth]

  % ---------- 正视图 ----------
  \draw (-2,6) -- (2,6) -- (1,5) -- (1,1) -- (2,0) -- (-2,0) -- (-1,1) -- (-1,5) -- cycle;
  \draw (-1,5) -- (1,5);
  \draw (-1,3) -- (1,3);
  \draw (-1,1) -- (1,1);

  % 正视图尺寸标注
  \draw[<->] (-2,6.6) -- (2,6.6);
  \draw (-2,6.45) -- (-2,6.75);
  \draw (2,6.45) -- (2,6.75);
  \node at (0,7.15) {$4$};
  \draw[<->] (-2,-0.6) -- (2,-0.6);
  \draw (-2,-0.75) -- (-2,-0.45);
  \draw (2,-0.75) -- (2,-0.45);
  \node at (0,-1.25) {$4$};
  \draw[<->] (2.75,6) -- (2.75,5);
  \node[right] at (2.85,5.5) {$1$};
  \draw[<->] (2.75,5) -- (2.75,3);
  \node[right] at (2.85,4) {$2$};
  \draw[<->] (2.75,3) -- (2.75,1);
  \node[right] at (2.85,2) {$2$};
  \draw[<->] (2.75,1) -- (2.75,0);
  \node[right] at (2.85,0.5) {$1$};
  \foreach \y in {0,1,3,5,6} {
    \draw (2.6,\y) -- (2.9,\y);
  }
  \draw[<->] (-1,4.2) -- (1,4.2);
  \node[above] at (0,4.2) {$2$};
  \node at (0,-2.1) {正视图};

  % ---------- 侧视图 ----------
  \begin{scope}[shift={(7,0)}]
    \draw (-2,6) -- (2,6) -- (1,5) -- (1,1) -- (2,0) -- (-2,0) -- (-1,1) -- (-1,5) -- cycle;
    \draw (-1,5) -- (1,5);
    \draw (-1,3) -- (1,3);
    \draw (-1,1) -- (1,1);
    \node at (0,-2.1) {侧视图};
  \end{scope}

  % ---------- 俯视图 ----------
  \draw (-2,-4) -- (2,-4) -- (2,-8) -- (-2,-8) -- cycle;
  \draw (0,-6) circle (2);
  \draw[dashed] (-1,-5) -- (1,-5) -- (1,-7) -- (-1,-7) -- cycle;
  \draw[dashed] (0,-6) circle (1);
  \draw (2,-4) -- (1,-5);
  \draw (2,-8) -- (1,-7);
  \draw (-2,-4) -- (-1,-5);
  \draw (-2,-8) -- (-1,-7);
  \node at (0,-9.1) {俯视图};

\end{tikzpicture}
